%% kapa: exam=1 round=21 year=2010 period=1 form=A subject=civil_law subject_ko=민법 seq=10 number=10
%% source: Q-Net 감정평가사 기출문제 (cst00309) artlSeq=5024181 file=extracted/1008/5024181/2010년 제21회 감정평가사 1차시험 기출문제/1교시.hwp
%% conversion: kordoc 4.24.51 -> tools/gen1.py (auto); review=proofread-sonnet; flags=none
\begin{question}{10}
부동산 물권변동에 관한 설명으로 옳지 않은 것은? (다툼이 있으면 판례에 의함)
\begin{choices}
  \item 건축물 도급계약에서 도급인과 수급인 사이에 완성된 건물의 소유권을 도급인에게 귀속시키기로 하는 합의가 있었다면, 그 후에 수급인이 자기의 노력과 재료를 들여 건물을 완성하였더라도 그 건물의 소유권은 도급인에게 원시적으로 귀속한다.
  \item 공유물분할판결에 의해 물권변동의 효력이 생기기 위해서는 등기를 필요로 한다.
  \item 부동산소유권이전등기절차의 이행을 구하는 소송에서 원고의 승소판결이 확정되더라도 그에 따른 이전등기를 하여야 물권변동의 효력이 생긴다.
  \item 구분소유적 공유관계의 공유지분에 설정된 근저당권의 실행으로 공유지분을 취득한 매수인은 그 구분소유적 공유지분을 그대로 승계취득한다.
  \item 혼동에 의한 물권의 소멸은 등기를 하지 않더라도 효력이 발생한다.
\end{choices}
\end{question}
\answer{2}
